\item \label{itm:rifting} {\bf [5 pts.]} Rifting \\
This question essentially comes from the start of the practical.

A simple model of rifting is illustrated in Figure~\ref{fig:rifting}.
\begin{figure}[hbt]
    \centering
    \includegraphics[]{figures/rifting.eps}
    \caption{Temperature}
    \label{fig:rifting}
\end{figure}
\begin{enumerate}
    \item Assuming the rifting takes place instantaneously, what are the initial
    and equilibrium geotherms (temperature as a function of depth)?

    \item As discussed in the practical and notes, the Fourier coefficients for
    this type of problem are given by
    \begin{equation}
        a_n = \frac{2}{L} 
        \int_0^L [f(z) - T_E(z)] \sin \left( \frac{n \pi z}{L} \right) \dd{z} .
    \end{equation}
    Using the geotherms determined above, what are the values of the Fourier
    coefficients?

    \item There are two timescales for this problem: one as cooling progresses
    from the base of the rifted lithosphere to the surface; and a second as
    cooling progresses from the base of the rifted lithosphere to the base of
    the equilibrium lithosphere.  What are these timescales?

    \item What happens as $\gamma \to 1$?  How does the timescale compare with
    this compare with that estimate by plate cooling?  Show your work.
\end{enumerate}
